% ============================================================
%  第三讲  三价构形上的拉格朗日 Floer 理论
% ============================================================
\section{第三讲：三价构形上的拉格朗日 Floer 理论}

本讲与 A. Efimov、L. Katzarkov 合作。目标是推广通常的 wrapped Floer 理论，允许非正合的拉格朗日，用 Novikov 系数记录全纯多边形的拓扑能量，再加上局部系统。

\subsection{纲领与柱面}
\origpage{11}

回忆裤子的 HMS：
\[
\{\text{裤子}\}\ \longleftrightarrow\ \bigl(\C^3,\,-z_1z_2z_3\bigr).
\]
前两讲用的是 U 形拉格朗日。现在希望改为在 $\crit(W)\subset Y$ 中工作，即工作在 $\C^*$ 与 $\PP^1$ 的三价构形上。

先把已知的对应摆出来。摩天层 $\OO_{(x_1,x_2)\in H}$ 对应坐标轴 $S^1\subset\C$ 上的顶针加界定余链（Aganagic--Vafa）；例如
\[
\val(x_1)>0\ \longleftrightarrow\ S^1\ \text{围出的面积},\qquad x_2=-1-x_1 ,
\]
其中 $S^1\subset\C$ 是拉格朗日。对摩天层 $\OO_H$ 本身，并不存在真正的顶针，但有一个分片光滑的组合，见图~\ref{fig:thicken}。

\begin{figure}[htbp]
\centering
\begin{tikzpicture}[scale=1.05,>=Stealth]
  % 三条坐标轴
  \draw[thick] (0,0) -- (150:1.5);
  \draw[thick] (0,0) -- (30:1.5);
  \draw[thick] (0,0) -- (270:1.5);
  \fill (0,0) circle (1.6pt);
  \node[below left] at (0,0) {\footnotesize $0$};
  % L_0 = 三条射线
  \draw[very thick,blue!70!black] (0,0) -- (150:1.05);
  \draw[very thick,blue!70!black] (0,0) -- (30:1.05);
  \draw[very thick,blue!70!black] (0,0) -- (270:1.05);
  \node[blue!70!black] at (215:1.35) {\footnotesize $L_0$};
  % L_p：一条腿上的圆
  \draw[thick,green!45!black] (30:0.62) ellipse (0.10 and 0.20);
  \node[green!45!black] at (30:0.95) {\footnotesize $L_p$};
  % phi(L_0)
  \draw[thick,red!70!black] (150:1.25) .. controls (90:1.05) and (60:0.85) .. (30:0.95);
  \node[red!70!black] at (95:1.35) {\footnotesize $\phi(L_0)$};
  \node at (0,-1.9) {\footnotesize $\Sigma=\cup_3\C$，$L_0=\cup_3\R_{\ge0}$};
\end{tikzpicture}
\caption{三价构形 $\Sigma=\cup_3\C$ 上的 $L_0=\cup_3\R_{\ge0}$、一条腿上的 $L_p\simeq S^1$，以及扰动后的 $\phi(L_0)$。（原稿第 11 页）}
\label{fig:sigma}
\end{figure}

\paragraph{柱面上的模型.} 先回忆柱面 $\R\times S^1$ 上的 wrapped Floer 同调，后面所有计算都以它为模板。
\begin{itemize}
  \item $\End(L_0)\simeq K\bigl[x^{\pm1}\bigr]$。
  \item 乘法 $x^i\cdot x^j=x^{i+j}$ 由三角形计数给出：三条边分别落在 $L_0$、$\phi(L_0)$、$\phi^2(L_0)$ 上。
  \item 在点处取值（evaluation）对应具体的复合 $L_0\to L_0\to L_p$。$L_p$ 上的边绕 $S^1$ 转 $k$ 圈（绕数 $k$），对应 $x^k$。
\end{itemize}

\paragraph{裤子上的对应.} 现在取 $\Sigma=\cup_3\C$，即"裤子的镜像"。考虑
\[
L_0=\cup_3\R_{\ge0}
\]
以及它的"wrapped Floer 同调"。要记住的是：这个对象代表热带拉格朗日那一部分，即 $\left(\C^*\right)^2$ 中 $\cap U\ell_0$ 的 $HW^0$。

计算给出：
\begin{itemize}
  \item $CW(L_0,\phi(L_0))$ 在三个柱形端各有无穷多个生成元；
  \item 顶点处还剩一个待定的生成元，结果恰好只有一个。
\end{itemize}
这些生成元应当给出 $\PP^1\setminus\{0,-1,\infty\}$（即裤子）上函数环的一组加法基。
\begin{itemize}
  \item 在点处取值：$L_p\simeq S^1$ 位于一条腿上，再加上局部系统，对应 $\OO_p$，其中 $p$ 是裤子上的点，靠近某个尖点。$S^1$ 在 $\C$ 中围出的面积对应局部坐标在该尖点处的赋值。
\end{itemize}

\begin{figure}[htbp]
\centering
\begin{tikzpicture}[scale=0.95,>=Stealth]
  \draw[thick] (0,0) -- (2.6,0) -- (2.6,0.5) -- (0,0.5) -- cycle;
  \draw[thick] (0,0.25) -- (2.6,0.25);
  \draw[thick] (0.9,-0.35) -- (0.9,0.85);
  \foreach \y in {0.08,0.42,0.70}{
    \draw[blue!70!black,thick] (0.05,\y) .. controls (0.7,\y+0.10) and (1.6,\y-0.08) .. (2.55,\y+0.05);
  }
  \node[above] at (0.9,0.9) {\footnotesize $L_0$};
  \node at (1.3,-0.45) {\footnotesize 柱面 $\R\times S^1$};
  \node[blue!70!black] at (2.05,0.95) {\footnotesize 三角形计数的边};
\end{tikzpicture}
\caption{柱面上乘法 $x^i\cdot x^j=x^{i+j}$ 的三角形：两条边落在 $L_0$ 上，第三条落在扰动后的像上。（原稿第 11 页）}
\label{fig:cylinder}
\end{figure}

\subsection{局部坐标、首项与传播}
\origpage{12}

$L_p$ 决定 $\FF(Y,w_0)$ 中对象的确切方式，依赖于界定余链的构造，也就是 Ag-Vafa 与 Lin 等人语言中的 framing。这归结为在裤子某一端选取局部坐标。

例如在裤子 $\{1+x_1+x_2=0\}\subset\left(\C^*\right)^2$ 中，靠近 $x_1\to0$ 处可以取 $x_1$，也可以取
\[
-\frac{x_1}{x_2}=\frac{x_1}{1+x_1},
\]
这两种选取最自然。注意
\[
\val(x_1)=\val\Bigl(-\frac{x_1}{x_2}\Bigr),
\]
两者都等于 $S^1$ 围出的面积。

\paragraph{首项.} 看首项，即最小的那些多边形，也就是完全落在柱形端内的多边形。无论对 $CW(L_0,L_0)$ 的试探生成元的乘积，还是在点处取值，这些首项的计算都与柱面情形一样。于是可以断定：柱形端内的生成元是局部坐标之逆的各次幂，也就是极点阶数递增的函数。验算如下：$1$ 与各端局部坐标的 $\pm k$ 次幂（$k\ge1$），三个端各一组，合起来构成 $\OO(H)$ 的一组加法基。这本质上就是部分分式分解。

\paragraph{传播.} 不过这些函数在裤子的其他端并不为零。因此必须有某些全纯圆盘穿过结点，把它们"传播"过去。

传播规则可以由 HMS 反推确定：预期它们与 $(Y,w_0)$ 中（或对热带拉格朗日）的计算相符。事后看，这些规则确实给出等价的 Fukaya 范畴，但目前没有直接证明。

% [Tutor Note: 易错] 光滑化参数是 $T^A$ 而不是 $T^{A/2}$，但半粘合的两式右端都是 $T^{A/2}$：一条 $\PP^1$ 的面积为 $A$，而它在每个结点处只贡献一半。写错这个指数，$p_1p_2p_3=1$ 就不再成立。来源：原书第 13 页 $yy'=T^A$ 与 $xy'=T^{A/2}$ 两式。
\subsection{\texorpdfstring{$A_{\mathrm{model}}$ 数据、B 侧解释与 $F$ 的对象}{A-model 数据、B 侧解释与 F 的对象}}
\origpage{13}

\paragraph{数据.} 三价构形的 $A_{\mathrm{model}}$ 数据包含两部分：
\begin{itemize}
  \item 各 $\PP^1$ 的辛面积 $A_1,A_2,A_3$；
  \item 每个结点处的"图卡数据"，它等价于裤子各端的一个局部坐标。
\end{itemize}
具体地，图卡数据是一组函数
\[
p_1,p_2,p_3\in\OO\bigl(\text{裤子}\cong\PP^1\setminus\{-1,0,\infty\}\bigr),
\]
要求 $p_1,p_2,p_3$ 都能延拓为 $\PP^1$ 上的亚纯函数，且仅在相应的尖点处有单极点。首选取法是让它们在另两个尖点处取值 $-1$ 与 $0$。例如：
\[
\begin{array}{lll}
\text{尖点 }\infty: & z \ \text{ 或 } -(z+1), & \\[2pt]
\text{尖点 }0: & -\dfrac{z+1}{z}\ \text{ 或 } \dfrac{1}{z}, & \\[6pt]
\text{尖点 }-1: & -\dfrac{1}{z+1}\ \text{ 或 } -\dfrac{z}{z+1}. &
\end{array}
\]

\paragraph{B 侧解释.} 在极大退化附近，光滑曲线是把各 $\PP^1$ 在 $\{-1,0,\infty\}$ 处粘起来得到的；每个结点处的光滑化形如
\[
y\,y'=T^{A},
\]
其中 $y,y'$ 是在标记点附近选取的局部坐标，取上面那些函数的逆；$T^{A}$ 是光滑化参数，$A$ 为该 $\PP^1$ 的面积。

对我们更方便的说法是把每次粘合拆成两半：
\[
\text{把 }\PP^1\ni0\text{ 粘到裤子上：}\quad x\,y'=T^{A/2};
\qquad
\text{把 }\PP^1\ni\infty\text{ 粘到裤子上：}\quad x'\,y=T^{A/2},
\]
其中 $x'=1/x$ 是另一个局部坐标。

图卡数据的选取反映了一个事实：若不选 $\PP^1$ 在结点附近的局部坐标，就没有完全典范的方式把裤子显式地粘起来。

\paragraph{传播规则.} 图卡数据决定 Floer 微分的传播规则。最典范的选取是
\[
\Bigl\{\,z,\ -\frac{z+1}{z},\ -\frac{1}{z+1}\,\Bigr\}
\qquad\text{或}\qquad
\Bigl\{\,-(z+1),\ \frac{1}{z},\ -\frac{z}{z+1}\,\Bigr\},
\]
它们显示循环对称性：只需指定结点处三条腿的循环次序即可。此时三个函数满足
\[
p_1p_2p_3=1,\qquad p_i^{-1}+p_{i+1}+1=0 .
\]

\begin{figure}[htbp]
\centering
\begin{tikzpicture}[scale=0.95,>=Stealth]
  % 左：裤子
  \draw[thick] (-1.6,0) ellipse (0.28 and 0.55);
  \draw[thick] (-1.6,0.55) .. controls (-0.2,0.75) and (0.2,0.75) .. (0.9,0.25);
  \draw[thick] (-1.6,-0.55) .. controls (-0.2,-0.75) and (0.2,-0.75) .. (0.9,-0.25);
  \draw[thick] (0.9,0.25) -- (1.55,0.62);
  \draw[thick] (0.9,-0.25) -- (1.55,-0.62);
  \node at (-1.6,0) {\footnotesize 裤子};
  \node at (-1.9,0.95) {\footnotesize $y$};
  % 结点
  \fill (1.0,0) circle (1.8pt);
  % 右：P^1
  \draw[thick] (1.05,0) ellipse (1.15 and 0.72);
  \node at (2.0,0.42) {\footnotesize $\PP^1$};
  \node at (1.35,0) {\footnotesize $x$};
  \node[below] at (1.6,-0.95) {\footnotesize $x\,y'=T^{A/2}$};
\end{tikzpicture}
\caption{B 侧：把 $\PP^1$ 在结点处粘到裤子上，光滑化参数 $T^{A/2}$。（原稿第 13 页）}
\label{fig:gluing}
\end{figure}

\paragraph{$F$ 的对象.} $F$ 中"线丛型"的对象是三价图，每个分支上有一条弧，并在结点处有固定方向；弧上还带局部系统，在结点处扭结。例如图~\ref{fig:objs} 所示的两种图形。

\begin{figure}[htbp]
\centering
\begin{tikzpicture}[scale=0.95,>=Stealth]
  % 第一幅：一条弧穿过一个顶点
  \begin{scope}
    \draw[thick] (0,0) circle (0.75);
    \draw[thick] (1.5,0) circle (0.75);
    \fill (0.75,0) circle (1.8pt);
    \draw[red!75!black,thick] (0,-0.75) arc (-90:90:0.75);
    \draw[red!75!black,thick] (1.5,-0.3) -- (0.85,-0.05);
    \draw[red!75!black,thick] (0.65,0.05) -- (1.5,0.3);
    \node at (0.75,-1.05) {\footnotesize (a)};
  \end{scope}
  \begin{scope}[xshift=3.6cm]
    \draw[thick] (0,0) circle (0.75);
    \draw[thick] (1.5,0) circle (0.75);
    \fill (0.75,0) circle (1.8pt);
    \draw[red!75!black,thick] (0.35,0.0) -- (1.15,0.0);
    \draw[red!75!black,thick] (0.75,0.35) -- (0.75,0.75);
    \node at (0.75,-1.05) {\footnotesize (b)};
  \end{scope}
\end{tikzpicture}
\caption{$\FF$ 的对象：带一条弧与局部系统的三价图（结点固定方向）。（原稿第 13 页）}
\label{fig:objs}
\end{figure}

为定义 $\hom(A,B)$，把 $A$ 在顶点处的腿略作逆时针推挤，并在任一分量中正向卷绕。

\subsection{生成元与线丛}
\origpage{14}

按上面的约定，每个顶点给出一个 $0$ 次生成元，其余交点处照常给出生成元。

\paragraph{解释.} 把线丛想成粘合而成：
\begin{itemize}
  \item 每条主干上的弧对应裤子上取 $\OO$；
  \item 每条 $\PP^1$ 上的弧绕 $k$ 圈，对应 $\PP^1$ 上取 $\OO(k)$；这里 $\PP^1$ 在 $0$ 与 $\infty$ 处有区分的三价点。
\end{itemize}

$\hom(L_1,L_2)$ 的生成元与下列对象对应：
\begin{itemize}
  \item 裤子上的生成元对应裤子上的 contact fusion；
  \item 在绕 $k$ 圈的 $\PP^1$ 上，有 $|k-1|$ 个生成元，对应该 $\PP^1$ 上在 $0$ 与 $\infty$ 处为零的 $\OO(k)$ 的截面；更准确地说，对应
  \[
  H^\bullet\Bigl(\PP^1,\ \LL_1^\vee\otimes\LL_2\otimes\OO(-p_0-p_\infty)\Bigr).
  \]
\end{itemize}

两种情形的行为不同：
\begin{itemize}
  \item 当 $k>0$ 时，端点处的生成元对应在 $0$ 与 $\infty$ 处不为零的单项式，即 $x^k$ 与 $x^{-k}$；
  \item 当 $k\le0$ 时，会出现 Floer 微分
  \[
  \partial a=T^{\mathrm{area}}\,b,\qquad \partial a'=T^{\mathrm{area}}\,b',
  \]
  使得复形的这一部分无论如何都有上同调 $\simeq H^\bullet\bigl(\PP^1,\OO(k)\bigr)$。
\end{itemize}

\paragraph{捏扁极限.} 若所有 $\PP^1$ 都被捏扁（面积为 $\infty$），等价地，B 侧不做光滑化，则只出现上面这些"细"的 Floer 微分，此时
\[
H^\bullet\hom(L_1,L_2)\ \simeq\ H^\bullet\bigl(C_0,\ \LL_1^\vee\otimes\LL_2\bigr),
\]
其中 $C_0$ 是未光滑化的结点曲线。$\hom$ 的生成元对应每个 $\PP^1$ 中 $y$ 坐标的单项式。

给 $\PP^1$ 有限面积会使 Floer 理论发生形变，允许比"细"曲线更多的全纯曲线；这正对应于把 $C_0$ 形变到上一小节描述的光滑化。

\begin{remark}
"捏扁"这个极限是整讲的算术核心。面积取 $\infty$ 时 Novikov 参数 $T^{\mathrm{area}}$ 变成 $0$，除端点外的高次项全部消失，微分只剩下"细带"部分；于是 Floer 复形逐项等于 $\PP^1$ 上的截面复形，同调由维数公式直接读出。面积有限时，$T^{\mathrm{area}}$ 不为零，高阶修正出现，$\hom$ 才真正依赖镜像曲线的模。
% [Tutor Note: 直觉] 面积 $\to\infty$ 相当于把 Novikov 参数推到 $0$。此时只剩最小的全纯曲线，同调退化为经典的上同调；有限面积时的全部困难都藏在"更大的全纯曲线"里。来源：原书第 14 页"捏扁"一段。
\end{remark}

\subsection{Floer 微分的提法}
\origpage{15}

现在给出 $\hom(L_1,L_2)="CF(L_1,L_2)"$ 上 Floer 微分的具体规则。做法是计数"刚性"的全纯带；所谓刚性指 $0$ 维，因而有限。这些带边界在 $L_1,L_2$ 上，把 $0$ 次生成元接到 $1$ 次生成元；它们可以穿过顶点，并按图卡数据的规定行事。

举一个例子：一条可能对 $\hom(A,B)$ 的 $\partial$ 有贡献的全纯带，输入为 $a$、输出为 $b$。它进入顶点时是"$1$ 次带"，也就是几乎一次覆盖圆盘；离开顶点时是"$0$ 次带"。

规则如下：
\begin{itemize}
  \item 设一条带以次数 $d_1\ge1$ 从局部坐标 $p_1$ 所标的腿进入顶点，以次数 $d_2\ge0$ 从标 $p_2$ 的腿离开。则这条带被允许，当且仅当在 $p_1^{d_1}$ 于 $p_2$ 的极点附近的幂级数展开中，$p_2^{-d_2}$ 的系数非零；此时它按该系数作为重数计数。
  \item 一条合规则的带，总权重为
  \[
  \pm\,T^{\mathrm{area}}\cdot\bigl(\text{边界上局部系统的 holonomy}\bigr)
  \cdot\prod\bigl(\text{带穿过的各顶点处的重数}\bigr).
  \]
\end{itemize}

\begin{example}\label{ex:rule}
取
\[
p_1=z,\qquad p_2=\frac{1}{z},\qquad p_3=-\frac{1}{z+1}.
\]
设带以次数 $d_1$ 从腿 $p_1$ 进入。它可以有两种出路：
\begin{enumerate}[label=(\arabic*)]
  \item 以次数 $d_2=d_1$ 从腿 $p_2$ 离开；
  \item 以任一 $0\le d_3\le d_1$ 从腿 $p_3$ 离开，离开时带因子 $\pm\binom{d_1}{d_3}$。
\end{enumerate}
理由是展开式
\[
p_1=\frac{1}{p_2}\ \Longrightarrow\ p_1^{d_1}=p_2^{-d_1},
\]
\[
p_1=-1-\frac{1}{p_3}\ \Longrightarrow\
p_1^{d_1}=(-1)^{d_1}\Bigl(1+d_1p_3^{-1}+\binom{d_1}{2}p_3^{-2}+\cdots+p_3^{-d_1}\Bigr).
\]
反过来，以次数 $d_2$ 从腿 $p_2$ 进入的带，既可以以次数 $d_1=d_2$ 从腿 $p_1$ 离开，也可以以任一 $d_3$ 从腿 $p_3$ 离开；这是因为
\[
p_2=-\frac{p_3}{p_3+1}=-1-p_3^{-1}-p_3^{-2}-\cdots .
\]
\end{example}

\begin{figure}[htbp]
\centering
\begin{tikzpicture}[scale=1.0,>=Stealth]
  % 三条腿
  \draw[thick] (0,0) -- (0,1.5);
  \draw[thick] (0,0) -- (-1.3,-0.75);
  \draw[thick] (0,0) -- (1.3,-0.75);
  \fill (0,0) circle (1.8pt);
  \node[above] at (0,1.5) {\footnotesize $p_2$};
  \node[below left] at (-1.3,-0.75) {\footnotesize $p_1$};
  \node[below right] at (1.3,-0.75) {\footnotesize $p_3$};
  % 带
  \draw[red!75!black,thick] (-1.05,-0.9) -- (-0.1,-0.1) -- (1.0,-0.75);
  \fill[red!75!black] (-1.05,-0.9) circle (1.4pt);
  \fill[red!75!black] (1.0,-0.75) circle (1.4pt);
  \node[red!75!black,below left] at (-1.05,-0.9) {\footnotesize $a$（输入，次数 $d_1$）};
  \node[red!75!black,below right] at (1.0,-0.75) {\footnotesize $b$（输出，次数 $d_3$）};
  \node at (0,-1.9) {\footnotesize 一条穿过顶点的全纯带};
\end{tikzpicture}
\caption{穿过顶点的全纯带：以次数 $d_1$ 从腿 $p_1$ 进入，以次数 $d_3$ 从腿 $p_3$ 离开，重数由 $p_1^{d_1}$ 的展开系数给出。（原稿第 15 页）}
\label{fig:strip}
\end{figure}

\subsection{裤子上的显式计算}
\origpage{16}

把上面的规则用到裤子上，可以算出 $CW(L_0,L_0)$。几条腿上的生成元 $x_0,x_1,x_2,\dots$ 与 $y$ 已在图~\ref{fig:cylinderend} 中标出。

\[
\hom(L_0,L_0)\ \longleftrightarrow\ H^\bullet\bigl(\text{亏格 }1\text{ 的曲线},\ \OO\bigr).
\]
微分作用如下：
\begin{itemize}
  \item $\partial x_0=0$：两条细带 $x_0\to y$ 相互抵消；
  \item $\partial x_1=\pm T^{\mathrm{area}}y$：唯一的轨迹如图，在 $z=0$ 附近
  \[
  -\frac{1}{z+1}=(-1)+\Bigl(\frac{1}{z}\Bigr)^{-1}-\Bigl(\frac{1}{z}\Bigr)^{-2}-\cdots,
  \]
  常数项为 $-1$，故允许在 $1/z$ 那条腿上以次数 $0$ 离开；
  \item $\partial x_2=\pm T^{\mathrm{something}}y$：同理，只有一条轨迹，进入次数 $2$、离开次数 $0$。
\end{itemize}
此外，在 $z=\infty$ 附近
\[
-\frac{1}{z+1}=-z^{-1}+z^{-2}-z^{-3}+z^{-4}-\cdots
\]
没有常数项。因此若从 $z$ 那条腿进入，带的次数必然 $\ge1$，随后就一直绕圈传播，永远无法终止在 $y$。

\begin{figure}[htbp]
\centering
\begin{tikzpicture}[scale=0.95,>=Stealth]
  % 裤子的三价图
  \draw[thick] (0,0) circle (0.9);
  \draw[thick] (-0.62,-0.65) -- (-1.5,-1.1);
  \draw[thick] (0.62,-0.65) -- (1.5,-1.1);
  \draw[thick] (0,0.9) -- (0,1.7);
  \draw[thick] (-0.2,-1.35) -- (-1.5,-1.35) -- (-1.5,-0.9);
  \fill (0,0) circle (1.8pt);
  \node at (0.45,0.45) {\footnotesize $x_0$};
  \node at (-0.95,0.35) {\footnotesize $x_1$};
  \node at (0.85,-0.25) {\footnotesize $x_2$};
  \node at (1.6,-0.75) {\footnotesize $y$};
  \node at (0,-2.1) {\footnotesize $\hom(L_0,L_0)$ 的生成元};
\end{tikzpicture}
\caption{裤子上 $CW(L_0,L_0)$ 的生成元分布（示意）。（原稿第 16 页）}
\label{fig:cylinderend}
\end{figure}

\paragraph{不对称性.} 上面两种走向的不对称，正是最终得到 $H^1\hom=0$ 的原因，这与预期相符。这个结果与"裤子上的函数、极点阶数为 $1$"一致。不过还存在另一组基，对每个极点阶数 $\ge2$ 各有一个生成元。事实上连环结构也能对上：作变量替换后得到
\[
K[x,y]\big/\bigl(y^2-x^3-ax-b\bigr),
\]
其中 $x\sim x_2+\cdots+x_1$，$y\sim x_3+\cdots$。这正是亏格 1 曲线的函数环。

\begin{remark}
这里是第三讲最漂亮的地方：一个纯辛几何的计算（数全纯带）自动产生了椭圆曲线的 $j$-不变量所编码的环结构。系数 $a,b$ 来自各 $\PP^1$ 的面积，也就是镜像曲线的模；而生成元的组合方式 $x\sim x_2+\cdots+x_1$ 恰好重现了 B 侧坐标在尖点处的展开。
% [Tutor Note: 直觉] 环结构不是额外加上去的，它是 $\mu^2$ 的直接产物；$a,b$ 由面积参数决定，把 $A$ 侧的面积数据翻译成了 $B$ 侧的模。来源：原书第 16 页环结构一段。
\end{remark}

\paragraph{高阶运算.} $\mu^2$ 及更高阶的运算按同样的方式定义：计数边界在适当扰动过的三价拉格朗日图上的全纯多边形，并沿从输入到输出的路径传播。换句话说，是"在结点处彼此粘接的圆盘构成的树"。若输出落在结点上，需要单独处理。系数域上的符号约定按通常的 $A_\infty$ 约定取。

\begin{theo}[A.--Efimov--Katzarkov，进行中]\label{thm:aek}
按上述方式定义的三价 $\C$ 与 $\PP^1$ 构形的 Fukaya 范畴，导出等价于镜像曲线上的凝聚层范畴；该曲线由粘合 $(\PP^1,\text{三点})$ 并如约光滑化而构造。
\end{theo}

% [Tutor Note: 点拨] 传播的重数就是幂级数展开的系数，而不是 $1$：一条带在顶点处分叉成若干支，每支各得一个系数，乘起来才是这条带的权重。这也解释了为什么传播规则完全由"图卡数据"决定，与面积无关。来源：原书第 17 页两个重数公式。
\subsection{\texorpdfstring{$\mu^2$ 的传播}{mu\^2 的传播}}
\origpage{17}

最后画出 $\mu^2$ 的传播图，把重数公式写清楚。

从左到右看：第一个圆盘以 $L_0,L_1$ 为边界，第二个圆盘以 $L_1,L_2$ 为边界，两者在结点处相接，输出在上方。三层结点串起来时，输出在上、输入在下。

\paragraph{穿过一个结点.} 设进入时次数为 $1$ 的坐标是 $w_1$，输出处次数为 $0$ 的坐标是 $w_2$。则重数由展开
\[
w_1^{\deg w_1}=c_0+c_1w^{-1}+c_2w^{-2}+\cdots
\]
的常数项 $c_0$ 给出；对 $\deg w_2=0$ 就取这一项。

\paragraph{输出落在结点上.} 设结点处两侧的坐标分别为 $w$（次数 $r$）与 $z$（次数 $s$）。则重数由展开
\[
w^{r}z^{s}=(c)+\sum_{k\ge1}a_kw^{k}+\sum_{\ell\ge1}a_\ell z^{\ell}
\]
中的常数项 $(c)$ 给出。

\begin{figure}[htbp]
\centering
\begin{tikzpicture}[scale=0.9,>=Stealth]
  % 圆盘树
  \begin{scope}
    \draw[thick] (0,0) circle (0.7);
    \draw[thick] (1.9,0) circle (0.7);
    \draw[thick] (0.95,1.25) circle (0.7);
    \fill (0.95,0) circle (1.8pt);
    \fill (0.95,1.25) circle (1.8pt);
    \node[left] at (-0.8,0) {\footnotesize $L_0$};
    \node[right] at (2.7,0) {\footnotesize $L_2$};
    \node[above] at (0.95,2.0) {\footnotesize output};
    \node at (0.95,-1.0) {\footnotesize 圆盘树};
  \end{scope}
  \draw[->,thick] (3.4,0.4) -- (4.3,0.4);
  % 结点处的展开
  \begin{scope}[xshift=5.0cm]
    \draw[thick] (-0.9,0) -- (0.9,0);
    \fill (0,0) circle (1.8pt);
    \node[above] at (-0.55,0.05) {\footnotesize $w$};
    \node[above] at (0.55,0.05) {\footnotesize $z$};
    \node[below] at (0,-0.5) {\footnotesize $w^{r}z^{s}=(c)+\sum a_kw^{k}+\sum a_\ell z^{\ell}$};
  \end{scope}
\end{tikzpicture}
\caption{$\mu^2$ 的传播：两个圆盘在结点处相接；结点处的重数由 $w^rz^s$ 展开的常数项给出。（原稿第 17 页）}
\label{fig:mu2}
\end{figure}

以上是第一讲至第三讲的全部内容。附录 A 是同一作者 2019 年报告幻灯片的译本，用更初等的语言重述了这里的纲领；附录 B 是全书术语对照表。
